% One landscape overview: route-specific lanes share a single checks legend.
\begingroup\hyphenpenalty=10000\exhyphenpenalty=10000
\begin{tikzpicture}[x=1cm,y=1cm,font=\small\sffamily,text=ink,
 panel/.style={draw=ink!20,rounded corners=3pt,fill=wash,inner sep=7pt,align=center},
 flow/.style={-{Latex[length=2mm]},draw=ink!65,line width=.8pt},
 bus/.style={draw=ink!65,line width=.8pt},
 decision/.style={diamond,aspect=1.85,draw=amber!70,fill=amber!8,align=center,
   inner sep=3pt,font=\small\sffamily\bfseries},
 caption/.style={font=\small\sffamily,align=left},
 heading/.style={font=\small\sffamily\bfseries,text=teal},
 edge/.style={font=\scriptsize\sffamily\bfseries,fill=white,inner sep=2pt}]

% Input and preparation occupy one strip, before the route decisions.
\node[panel,minimum width=24.8cm,minimum height=1.65cm] (prep) at (12.5,0) {};
\pic[scale=.55] at (.45,.5) {dqscan};
\node[caption,anchor=west,text width=6.5cm] at (1.55,0) {
  \textbf{INPUT}\quad PDF or image\\Schema / template + optional task};
\node[caption,anchor=west,text width=15.5cm] at (8.3,0) {
  \textbf{1 / PREPARE}\quad Normalize PDF; count pages; inspect embedded text.\\
  Keep page images and text; request a MinerU parse for scans when available.};

\node[heading] at (1.95,-1.35) {2 / ROUTE};
\node[heading] at (7.2,-1.35) {3 / READ};
\node[heading] at (13,-1.35) {4 / CHECK};
\node[heading] at (18,-1.35) {5 / SELECT};
\node[heading] at (23.15,-1.35) {6 / RETURN};

% Decisions run downwards; their yes branches feed the corresponding lane.
\node[decision] (layoutq) at (1.95,-3) {Specialized\\layout?};
\node[decision] (textq) at (1.95,-5.5) {Enough\\PDF text?};
\node[decision] (shortq) at (1.95,-8) {Scan has\\$\leq12$ pages?};
\node[panel,fill=amber!8,minimum width=2.8cm,font=\small\sffamily\bfseries]
  (long) at (1.95,-10.5) {Long scan\\13+ pages};
\draw[flow] (.2,-.83)--(.2,-3)--(layoutq.west);
\draw[flow] (layoutq.south)--node[edge,right] {no} (textq.north);
\draw[flow] (textq.south)--node[edge,right] {no} (shortq.north);
\draw[flow] (shortq.south)--node[edge,right] {no} (long.north);

% Four document pictures distinguish the reading strategies.
\foreach \y/\name in {-3/layout,-5.5/text,-8/short,-10.5/longread} {
  \node[panel,fill=teal!4,minimum width=5.3cm,minimum height=2.15cm]
    (\name) at (7.2,\y) {};
}
\draw[flow] (layoutq.east)--node[edge,above] {yes} (layout.west);
\draw[flow] (textq.east)--node[edge,above] {yes} (text.west);
\draw[flow] (shortq.east)--node[edge,above] {yes} (short.west);
\draw[flow] (long.east)--(longread.west);

\pic[scale=.57] at (4.85,-2.47) {dqlayout};
\pic[scale=.57] at (4.85,-4.97) {dqtext};
\pic[scale=.57] at (4.85,-7.47) {dqscan};
\pic[scale=.57] at (4.85,-9.97) {dqscan};
\node[caption,anchor=west,text width=3.6cm] at (5.95,-3) {
  \textbf{A / Layout}\\Column crops or ruled bands\\\textcolor{teal}{One layout reading}};
\node[caption,anchor=west,text width=3.6cm] at (5.95,-5.5) {
  \textbf{B / Text PDF}\\Text + images when allowed\\\textcolor{teal}{One primary reading}};
\node[caption,anchor=west,text width=3.6cm] at (5.95,-8) {
  \textbf{C / Short scan}\\Primary + parse* + one-shot\\\textcolor{teal}{Up to three readings}};
\node[caption,anchor=west,text width=3.6cm] at (5.95,-10.5) {
  \textbf{D / Long scan}\\Primary, then parse*\\\textcolor{teal}{Up to two readings}};

% One shared stage, not a merge: every JSON gets its own source-specific chain.
\node[panel,minimum width=3.9cm,minimum height=9.65cm] (checks) at (13,-6.75) {};
\node[caption,align=center,text width=3.4cm] at (13,-2.8) {
  \textbf{Check each JSON}\\Use its source's chain.};
\draw[ink!18] (11.4,-3.5)--(14.6,-3.5);
\node[caption,align=center,text width=3.4cm] at (13,-4.45) {
  \textbf{Primary}\\Repairs $\to$\\output checks};
\node[caption,align=center,text width=3.4cm] at (13,-6.05) {
  \textbf{Parse}\\Spacing $\to$ repairs\\$\to$ output checks};
\node[caption,align=center,text width=3.4cm] at (13,-7.65) {
  \textbf{Layout / one-shot}\\Output checks};
\draw[ink!18] (11.4,-8.55)--(14.6,-8.55);
\pic at (13,-9.25) {dqcheck};
\node[caption,align=center,text width=3.3cm] at (13,-10.45) {
  Conditional repairs.\\No model calls.};

% Route-specific selection remains aligned with each input lane.
\foreach \y/\name in {-3/keepLayout,-5.5/keepText,-8/keepShort,-10.5/keepLong} {
  \node[panel,fill=amber!7,minimum width=4.3cm,minimum height=2.15cm]
    (\name) at (18,\y) {};
  \draw[flow] (9.85,\y)--(11.05,\y);
  \draw[flow] (14.95,\y)--(15.85,\y);
}
\node[caption,align=center,text width=3.8cm] at (18,-3) {\textbf{Keep layout reading}};
\node[caption,align=center,text width=3.8cm] at (18,-5.5) {\textbf{Keep primary reading}};
\node[caption,align=center,text width=3.8cm] at (18,-8) {
  \textbf{Successful readings}\\1: return it\\2--3: quorum\\0: error};
\node[caption,align=center,text width=3.8cm] at (18,-10.5) {
  \textbf{Prefer successful parse}\\Otherwise keep primary;\\error if neither succeeds.};

% Only the chosen successful route reaches the output.
\draw[bus] (20.75,-3)--(20.75,-10.5);
\foreach \y in {-3,-5.5,-8,-10.5} {
  \draw[bus] (20.15,\y)--(20.75,\y);
  \fill[ink!65] (20.75,\y) circle (1.2pt);
}
\draw[flow] (20.75,-6.75)--(21.5,-6.75);
\node[panel,fill=teal!7,minimum width=3.3cm,minimum height=6cm] at (23.15,-6.75) {};
\pic[scale=.7] at (22.65,-4.25) {dqjson};
\node[caption,align=center,text width=2.9cm] at (23.15,-6.5) {
  \textbf{Extracted JSON}\\Refill template\\if supplied.};
\draw[teal!25] (21.9,-7.6)--(24.4,-7.6);
\node[caption,align=center,text width=2.9cm] at (23.15,-8.55) {
  \textbf{Selection report}\\Readings used and\\selected candidate};

\node[anchor=north west,align=left,text width=24.6cm,font=\footnotesize\sffamily,text=ink!75]
  at (.15,-12) {
  \textbf{Only one route runs.} A reading is a candidate JSON. Quorum keeps one complete candidate; it does not mix fields.\\[4pt]
  * Parse reading requires cached pages. Cached text is used only at $\geq90\%$ page-count coverage; otherwise the reader uses PDF text.};
\end{tikzpicture}
\endgroup
